\documentclass[12pt]{article}
\usepackage{amsmath,amssymb,amsfonts}
\usepackage[margin=1in]{geometry}

\begin{document}

\textbf{Chapter 8, Problem 21.} At the beginning of Section 8.6 in discussing kernels of finite rank, we introduced briefly the notion of a linear functional. The purpose of this problem and the two that follow is to introduce some of the simpler aspects of linear functionals on a Banach space. We make the following definition: A \emph{functional} $F$ on a Banach space $B$ is an operation which assigns to every $g\in B$ a complex number (assuming that $B$ is defined over the field of complex numbers). The action of $F$ on $g$ is denoted by $F(g)$. $F$ is said to be \emph{linear} if $F(\alpha g_1+\beta g_2) = \alpha F(g_1)+\beta F(g_2)$ for all complex numbers $\alpha$ and $\beta$, and all $g_1,g_2\in B$. $F$ is said to be \emph{bounded} if there exists a real number $M$ such that $|F(g)|\le M\|g\|$ for all $g\in B$, $M$ being independent of $g$.

\textbf{(a)} Let $F$ be a bounded, linear functional on $B$. Suppose that we have in $B$ a sequence $\{g_n\}$ which tends in norm to $f\in B$. Show that $F(f_n) = F(f)$, that is, that $F$ is continuous.

\textbf{(b)} Show that if $f\in L_q$, then the functional $F$ defined on $L_p$ ($1/p+1/q=1$) by
\[
F(g) = \int_a^b f(x)\overline{g(x)}\,dx
\]
is a bounded, linear functional.

(\textit{Hint:} Look back at Problem 8.2.) It can be shown that every bounded, linear functional on $L_p$ can be written in this form, i.e., can be generated by an element of $L_q$.

\bigskip
\textbf{Solution.}

\textbf{Part (a):}
Let $g_n\to f$ in norm, i.e., $\|g_n-f\|\to 0$. Since $F$ is bounded and linear:
\[
|F(g_n)-F(f)| = |F(g_n-f)| \le M\|g_n-f\| \to 0.
\]
Therefore $F(g_n)\to F(f)$, proving continuity. $\blacksquare$

\textbf{Part (b):}

\emph{Linearity:} For $\alpha,\beta\in\mathbb{C}$ and $g_1,g_2\in L_p$:
\[
F(\alpha g_1+\beta g_2) = \int f\overline{(\alpha g_1+\beta g_2)}\,dx = \bar\alpha\int f\bar{g_1}\,dx + \bar\beta\int f\bar{g_2}\,dx = \bar\alpha F(g_1)+\bar\beta F(g_2).
\]

(Note: If the convention is $F(\alpha g) = \alpha F(g)$ rather than $\bar\alpha F(g)$, then the functional is conjugate-linear in the Hilbert space sense. With the standard Banach space definition of linearity $F(\alpha g) = \alpha F(g)$, one should instead define $F(g) = \int \bar{f}g\,dx$. We proceed with the convention that gives linearity.)

Using $F(g) = \int \bar{f}(x)g(x)\,dx$ for proper linearity:
\[
F(\alpha g_1+\beta g_2) = \alpha F(g_1)+\beta F(g_2). \quad\checkmark
\]

\emph{Boundedness:} By H\"older's inequality (Problem 8.2):
\[
|F(g)| = \left|\int \bar{f}\,g\,dx\right| \le \|f\|_q\|g\|_p.
\]
So $F$ is bounded with $M = \|f\|_q$:
\[
|F(g)| \le \|f\|_q\|g\|_p \quad\text{for all }g\in L_p. \quad\blacksquare
\]

\end{document}
